% GENERATED FILE. Edit canonical lessons, then run npm run generate:books.
% canonical-source-hash: fnv1a64:aa54e3e4bd409d12
% canonical-lessons: GU-C402-chakdo, GU-C402-tempo, GU-C402-gadu, GU-C402-kandaktar, GU-C402-hakvu

\chapter{Carts and Chhakdas}
\label{ch:gu-carts-and-chhakdas}
\hlchaptermodality{402}

\begin{chapteropening}
\textbf{By the end of this chapter:} \emph{I can talk about the carts and vans on a country road.}

Complete the last lesson of chapter 402: I can talk about the carts and vans on a country road.
\end{chapteropening}

\section[\texorpdfstring{\gu{છકડો} (chakḍo)}{છકડો (chakḍo)}]{\gu{છકડો} (chakḍo) — a chhakda (a country motor cart)}

\label{lesson:GU-C402-chakdo}



\begin{quote}
Before the new one: say the Gujarati for to drag, then the Gujarati for to overturn.
\end{quote}

\subsection*{You'll want to know: \gu{છકડો}}
\textbf{\gu{છકડો}} — \emph{chakḍo} — \textquotedblleft{}a chhakda (a country motor cart)\textquotedblright{}.

\begin{quote}
\gu{ગામના} \gu{માણસો} \gu{છકડામાં} \gu{બેસીને} \gu{શહેર} \gu{ગયા}.
\end{quote}

\begin{quote}
\emph{gāmnā māṇso chakḍāmā̃ besīne śaher gayā.}
\end{quote}

\begin{quote}
The village people rode into town in a chhakda.
\end{quote}

A masculine noun.

\begin{cousinweb}
\textbf{\gu{છકડો}} (\emph{chakḍo}) is a tadbhava of Sanskrit \emph{śakaṭa}, \textquotedblleft{}a cart\textquotedblright{}; CDIAL lists the Gujarati word for a cart (CDIAL 12236). Today the name belongs to the open three-wheeled motor cart of the Saurashtra roads.
\end{cousinweb}

\begin{grammarlens}[title={What you've built}]
One more word for getting around.
\end{grammarlens}

\subsection*{Your turn}
\begin{itemize}
\raggedright
  \item \emph{Say it:} \emph{chakḍo}
  \item \emph{Say it:} \emph{chakḍo}, once more
  \item \emph{Say it:} \emph{gāmnā māṇso chakḍāmā̃ besīne śaher gayā.} — The village people rode into town in a chhakda.
  \item \emph{Recall:} say the Gujarati for to drag, then the Gujarati for to overturn, then say \emph{chakḍo} again
\end{itemize}

\begin{tcolorbox}[breakable,colback=teal!4,colframe=teal!35!black,title={Before you move on}]
What does \gu{છકડો} mean? (\textquotedblleft{}A chhakda (a country motor cart)\textquotedblright{}.) Say it once more.
\end{tcolorbox}
\section[\texorpdfstring{\gu{ટેમ્પો} (ṭempo)}{ટેમ્પો (ṭempo)}]{\gu{ટેમ્પો} (ṭempo) — a tempo (a small goods van)}

\label{lesson:GU-C402-tempo}



\begin{quote}
Before the new one: say the Gujarati for to overturn, then the Gujarati for a chhakda.
\end{quote}

\subsection*{You'll want to know: \gu{ટેમ્પો}}
\textbf{\gu{ટેમ્પો}} — \emph{ṭempo} — \textquotedblleft{}a tempo (a small goods van)\textquotedblright{}.

\begin{quote}
\gu{અમે} \gu{ટેમ્પોમાં} \gu{બધો} \gu{સામાન} \gu{ભર્યો}.
\end{quote}

\begin{quote}
\emph{ame ṭempomā̃ badho sāmān bharyo.}
\end{quote}

\begin{quote}
We loaded all the luggage into a tempo.
\end{quote}

A masculine noun.

\begin{cousinweb}
\textbf{\gu{ટેમ્પો}} (\emph{ṭempo}) is borrowed from Indian English \emph{tempo}, which comes from the name of a firm whose three-wheeled goods vans were common in India.
\end{cousinweb}

\begin{grammarlens}[title={What you've built}]
One more word for getting around.
\end{grammarlens}

\subsection*{Your turn}
\begin{itemize}
\raggedright
  \item \emph{Say it:} \emph{ṭempo}
  \item \emph{Say it:} \emph{ṭempo}, once more
  \item \emph{Say it:} \emph{ame ṭempomā̃ badho sāmān bharyo.} — We loaded all the luggage into a tempo.
  \item \emph{Recall:} say the Gujarati for to overturn, then the Gujarati for a chhakda, then say \emph{ṭempo} again
\end{itemize}

\begin{tcolorbox}[breakable,colback=teal!4,colframe=teal!35!black,title={Before you move on}]
What does \gu{ટેમ્પો} mean? (\textquotedblleft{}A tempo (a small goods van)\textquotedblright{}.) Say it once more.
\end{tcolorbox}
\section[\texorpdfstring{\gu{ગાડું} (gāḍũ)}{ગાડું (gāḍũ)}]{\gu{ગાડું} (gāḍũ) — a cart (a bullock cart)}

\label{lesson:GU-C402-gadu}



\begin{quote}
Before the new one: say the Gujarati for a chhakda, then the Gujarati for a small goods van.
\end{quote}

\subsection*{You'll want to know: \gu{ગાડું}}
\textbf{\gu{ગાડું}} — \emph{gāḍũ} — \textquotedblleft{}a cart (a bullock cart)\textquotedblright{}.

\begin{quote}
\gu{બળદનું} \gu{ગાડું} \gu{ધીમે} \gu{ચાલે} \gu{છે}.
\end{quote}

\begin{quote}
\emph{baḷadnũ gāḍũ dhīme cāle che.}
\end{quote}

\begin{quote}
The bullock cart moves slowly.
\end{quote}

A neuter noun.

\begin{cousinweb}
\textbf{\gu{ગાડું}} (\emph{gāḍũ}) goes back to \emph{*gāḍḍa}, \textquotedblleft{}a cart\textquotedblright{}, a form CDIAL reconstructs; CDIAL lists the Gujarati word (CDIAL 4116). \gu{ગાડી}, a car, is its relative: the cart came first.
\end{cousinweb}

\begin{grammarlens}[title={What you've built}]
One more word for getting around.
\end{grammarlens}

\subsection*{Your turn}
\begin{itemize}
\raggedright
  \item \emph{Say it:} \emph{gāḍũ}
  \item \emph{Say it:} \emph{gāḍũ}, once more
  \item \emph{Say it:} \emph{baḷadnũ gāḍũ dhīme cāle che.} — The bullock cart moves slowly.
  \item \emph{Recall:} say the Gujarati for a chhakda, then the Gujarati for a small goods van, then say \emph{gāḍũ} again
\end{itemize}

\begin{tcolorbox}[breakable,colback=teal!4,colframe=teal!35!black,title={Before you move on}]
What does \gu{ગાડું} mean? (\textquotedblleft{}A cart (a bullock cart)\textquotedblright{}.) Say it once more.
\end{tcolorbox}
\section[\texorpdfstring{\gu{કંડક્ટર} (kanḍakṭar)}{કંડક્ટર (kanḍakṭar)}]{\gu{કંડક્ટર} (kanḍakṭar) — a bus conductor}

\label{lesson:GU-C402-kandaktar}



\begin{quote}
Before the new one: say the Gujarati for a small goods van, then the Gujarati for a bullock cart.
\end{quote}

\subsection*{You'll want to know: \gu{કંડક્ટર}}
\textbf{\gu{કંડક્ટર}} — \emph{kanḍakṭar} — \textquotedblleft{}a bus conductor\textquotedblright{}.

\begin{quote}
\gu{બસના} \gu{કંડક્ટરે} \gu{મને} \gu{ટિકિટ} \gu{આપી}.
\end{quote}

\begin{quote}
\emph{basnā kanḍakṭare mane ṭikiṭ āpī.}
\end{quote}

\begin{quote}
The bus conductor gave me a ticket.
\end{quote}

A masculine noun.

\begin{cousinweb}
\textbf{\gu{કંડક્ટર}} (\emph{kanḍakṭar}) is borrowed from English \emph{conductor}, the one who sells tickets on a bus.
\end{cousinweb}

\begin{grammarlens}[title={What you've built}]
One more word for getting around.
\end{grammarlens}

\subsection*{Your turn}
\begin{itemize}
\raggedright
  \item \emph{Say it:} \emph{kanḍakṭar}
  \item \emph{Say it:} \emph{kanḍakṭar}, once more
  \item \emph{Say it:} \emph{basnā kanḍakṭare mane ṭikiṭ āpī.} — The bus conductor gave me a ticket.
  \item \emph{Recall:} say the Gujarati for a small goods van, then the Gujarati for a bullock cart, then say \emph{kanḍakṭar} again
\end{itemize}

\begin{tcolorbox}[breakable,colback=teal!4,colframe=teal!35!black,title={Before you move on}]
What does \gu{કંડક્ટર} mean? (\textquotedblleft{}A bus conductor\textquotedblright{}.) Say it once more.
\end{tcolorbox}
\section[\texorpdfstring{\gu{હાંકવું} (hā̃kvũ)}{હાંકવું (hā̃kvũ)}]{\gu{હાંકવું} (hā̃kvũ) — to drive a cart; to drive (animals) on}

\label{lesson:GU-C402-hakvu}



\begin{quote}
Before the new one: say the Gujarati for a bullock cart, then the Gujarati for a bus conductor.
\end{quote}

\subsection*{You'll want to know: \gu{હાંકવું}}
\textbf{\gu{હાંકવું}} — \emph{hā̃kvũ} — \textquotedblleft{}to drive a cart; to drive (animals) on\textquotedblright{}.

\begin{quote}
\gu{દાદા} \gu{રોજ} \gu{સવારે} \gu{ગાડું} \gu{હાંકે} \gu{છે}.
\end{quote}

\begin{quote}
\emph{dādā roj savāre gāḍũ hā̃ke che.}
\end{quote}

\begin{quote}
Grandfather drives the cart every morning.
\end{quote}

Past, with \gu{તેણે} (he or she …): \gu{તેણે} \gu{ગાડું} \gu{હાંક્યું}, \emph{teṇe gāḍũ hā̃kyũ}.

\begin{cousinweb}
\textbf{\gu{હાંકવું}} (\emph{hā̃kvũ}) is a tadbhava of Sanskrit \emph{hakkayati}, \textquotedblleft{}calls to\textquotedblright{} (CDIAL 13939): you drive oxen on by shouting at them. CDIAL lists the Gujarati verb, \textquotedblleft{}to drive\textquotedblright{}.
\end{cousinweb}

\begin{grammarlens}[title={What you've built}]
One more word for getting around.
\end{grammarlens}

\subsection*{Your turn}
\begin{itemize}
\raggedright
  \item \emph{Say it:} \emph{hā̃kvũ}
  \item \emph{Say it:} \emph{hā̃kvũ}, once more
  \item \emph{Say it:} \emph{dādā roj savāre gāḍũ hā̃ke che.} — Grandfather drives the cart every morning.
  \item \emph{Recall:} say the Gujarati for a bullock cart, then the Gujarati for a bus conductor, then say \emph{hā̃kvũ} again
\end{itemize}

\begin{tcolorbox}[breakable,colback=teal!4,colframe=teal!35!black,title={Before you move on}]
What does \gu{હાંકવું} mean? (\textquotedblleft{}To drive a cart; to drive (animals) on\textquotedblright{}.) Say it once more.
\end{tcolorbox}
